\section{Change a representation without changing the interface}\label{rgb-purpose}
A pixel format specifies how component values occupy bits and bytes. This project
converts a sequence between two such formats. Start with a scalar conversion that
handles every legal count and byte alignment. Then replace groups of eight scalar
iterations with vector operations while retaining the same tail and public result.
The independent scalar implementation remains available as a comparison.

A smaller representation cannot retain every original color. Three eight-bit
components have $2^{24}$ possible combinations; sixteen packed bits have only
$2^{16}$. Several inputs must therefore select the same output. Establish which
bits are discarded before comparing code or judging a displayed result. A different
rounding rule may be useful in another application, but it would implement a
different conversion contract here.

The kernel operates on component code values. It does not apply a transfer function,
perform gamma correction, blend pixels or convert between color spaces. Those are
separate operations with their own arithmetic and display assumptions. The examples
below can be checked from bytes and integer formulas without relying on a monitor.
